%! Author = sriadityavedantam
%! Date = 3/25/26

\section{Some Preliminaries}

\lesson{1}{Wednesday 13 August 2025 09:10}{Day 1}

Though things arise naturally out of studies like Algebra, some claim that we need to go beyond
such, and apply some intuition and \('\)analytical\('\) prowess to find ideas from what may seem
abstract or uncorrelated.
Real Analysis gave foundation and birth to many different fields we
know today, including most of the applied math workforce, such as Chemists, Physicists,
Biologists, Meteorologists, and so forth.
But within math, also optimization, complex analysis,
differential equations, differential geometry, and so forth.
It is essentially a cornerstone of
mathematics, regarded by many.

\begin{definition}[Real Numbers]
	\textbf{The complete ordered} field that contains $\q$.
\end{definition}

\begin{definition}[Rational Numbers]
    \[\q = \{\dfrac{a}{b} : a,b\in\z, b\neq 0\}\]
\end{definition}

\begin{definition}[Integers]
	\[\z = \{\pm n : n\in\n\}\cup\{0\}\]
\end{definition}

\begin{definition}[Naturals]
	\[\n = \{1,2,\ldots\}\]
\end{definition}

\begin{example}[Another Defintion of Rationals]
    It can be said that the rationals can be defined like the following.
	\[\q = \z\times (\z\setminus\{0\})\]
\end{example}

\begin{remark}
    Note that when were asserting $\q$, then we were unable to use the defining operator : $\coloneqq$.
\end{remark}

\begin{corollary}{
    Suppose we have $(a,b), (\alpha,\beta)\in\z\times(\z\setminus\{0\})$. Then $(a,b)\sim(\alpha,\beta)$ if and only if $\alpha b = a\beta$.
}
\end{corollary}

\begin{definition}[Equivalence Relation]
    Equivalence relations have three properties.
    Reflexivity, Symmetry, and Transitivity.
	The reflexive property holds $(a,b)\sim(a,b)$ is true.
	The symmetric property holds if $(a,b)\sim(\alpha,\beta)$, then $(\alpha,\beta)\sim(a,b)$.
	The transitive property holds if $(a,b)\sim(\alpha, \beta)$ and $(\alpha,\beta)\sim(c,d)$, then $(a,b)\sim(c,d)$.
	Prove this property for Exercise 1.
\end{definition}

\begin{remark}
Hence, we can define the rationals as the following.
	\[\q \coloneqq \dfrac{\z\times(\z\setminus\{0\})}{\sim}\]
\end{remark}

\begin{definition}[Equivalence Class]
	\[[(a,b)] + [(\alpha,\beta)] = [(a\beta + \alpha b, b\beta)]\].
\end{definition}

\begin{proposition*}[Addition is well-defined.]{
	Given $(a,b)\sim(c,d)$ and $(\alpha,\beta)\sim(\gamma,\delta)$, then $(a \beta + \alpha b, b\beta)\sim(c \delta + \beta d, d\delta)$. Prove this for Exercise 2.
}
\end{proposition*}

\begin{lemma}[Irrationality Proof]{
	There is no $q\in\q$ such that $q^2 = 2$. We call this, $\sqrt{2}$ irrational.
}
	Suppose for the sake of contradiction, if $q = \dfrac{a}{b}$ $a,b\in\z$ and $a,b$ coprime, then $q^2 = \dfrac{a^2}{b^2} = 2$.
	Thus $a^2 = 2b^2$, then $a^2$ is even.
	Furthermore, $a = 2k$ for some $k\in\z$, therefore $a^2 = 4k^2$.
	Hence, $b^2$ is even, which contradicts our coprime assumption.
\end{lemma}

\begin{remark}
	The direct statement of this lemma is:
	\begin{displayquote}
		If $q^2 = 2$, then $q\notin\q$.
	\end{displayquote}
	A contrapositive of this lemma is:
	\begin{displayquote}
		If $q\in\q$, then $q^2\neq 2$.
	\end{displayquote}
\end{remark}